\section{Discussion \& Limitations}
\label{sec:discussion}

\textbf{Why does architecture search hurt on SMILES?}
We hypothesize that the SMILES domain is ``simple enough''---small vocabulary, short sequences, high regularity---that the default architecture is already near-optimal. Architectural modifications waste experimental budget on structural changes that do not improve over the default, while the HP-only agent efficiently tunes learning rates and schedules within a fixed, adequate architecture. This suggests a ceiling effect: when the architecture is not a bottleneck, searching over it is counterproductive.

\textbf{The clustering-universality paradox.}
Architectures cluster by domain ($p = 0.004$), yet all innovations transfer freely (41/41 universal). We interpret this as evidence that search dynamics---what the agent explores first, conditioned on early training signals---drive apparent specialization. The agent follows different optimization paths for SMILES and NLP data, but the structural innovations it discovers along those paths happen to be universally effective at the ${\sim}$10M parameter scale. Whether this universality holds at larger scales, where domain-specific patterns may become more important, remains an open question.

\textbf{Limitations.}
(1)~\emph{Small model scale} (${\sim}$8.6M parameters). At larger scales, domain-specific architectural features (e.g., attention patterns tuned to molecular bonding topology) may become more valuable, potentially changing the decomposition balance.
(2)~\emph{Short training proxy} (5 minutes). The calibration study yields $\rho = 0.54$, indicating moderate rank correlation with longer training. Architecture rankings could change with extended training budgets, particularly for architectures with different convergence rates.
(3)~\emph{Single LLM backend} (Claude Sonnet). The agent's search behavior---and the resulting decomposition---may differ with other LLMs or with different prompt engineering.
(4)~\emph{Low statistical power} ($n = 3$--$5$ runs per condition). The protein track in particular cannot distinguish between conditions. Larger-scale replication would strengthen the protein and NLP conclusions.
(5)~\emph{SMILES representation only}. SMILES is one of several molecular string representations; SELFIES, InChI, or 3D coordinate representations may yield different decomposition patterns.
(6)~\emph{Closed-source agent}. The LLM agent operates via API, introducing non-determinism. We mitigate this by running multiple independent replicates per condition, but exact reproduction of individual runs is not possible.
